\documentclass[12pt,a4paper]{article}

\usepackage[utf8]{inputenc}
\usepackage[T1]{fontenc}
\usepackage[brazil]{babel}
\usepackage{geometry}
\usepackage{setspace}

\geometry{left=3cm,right=2cm,top=3cm,bottom=2cm}
\onehalfspacing

\title{Análise da Evasão Acadêmica em Universidade Particular Utilizando Técnicas de Machine Learning}
\author{Guilherme Machado}
\date{2026}

\begin{document}

\maketitle

\begin{abstract}
Este estudo tem como objetivo analisar e prever a evasão acadêmica em uma universidade particular utilizando técnicas estatísticas e algoritmos de machine learning. Busca-se identificar fatores relevantes associados à evasão e propor ações preventivas.
\end{abstract}

\section{Introdução}
A evasão acadêmica é um problema recorrente em instituições de ensino superior, impactando a sustentabilidade financeira e o desenvolvimento educacional dos estudantes.

\section{Problema de Pesquisa}
Quais fatores influenciam a evasão acadêmica e como técnicas de machine learning podem ser utilizadas para prever esse fenômeno?

\section{Objetivos}

\subsection{Objetivo Geral}
Analisar e prever a evasão acadêmica utilizando técnicas de machine learning.

\subsection{Objetivos Específicos}
\begin{itemize}
\item Identificar variáveis relevantes
\item Aplicar modelos estatísticos
\item Desenvolver modelos preditivos
\item Comparar algoritmos
\end{itemize}

\section{Metodologia}
Pesquisa quantitativa com uso de dados acadêmicos, financeiros e demográficos. Aplicação de regressão logística, árvores de decisão e Random Forest.

\section{Resultados Esperados}
Identificação de fatores de risco e apoio à tomada de decisão institucional.

\section{Referências}
JAMES, G. et al. \textit{An Introduction to Statistical Learning}. Springer, 2013.

\end{document}